% problem_id: S001-spaces-pushouts-theorem-nagata
% source_id: S001 (Stacks Project)
% source_file: spaces-pushouts.tex
% statement_locator: spaces-pushouts.tex lines 3581--3589
% proof_locator: spaces-pushouts.tex lines 3591--3664
% env: theorem
% id_note: label 跨文件重名 ⇒ id 用文件限定形式
% status: complete (statement + author proof verbatim + reason + locators)
%
\begin{theorem}
\label{theorem-nagata}
\begin{reference}
\cite{CLO}
\end{reference}
Let $S$ be a scheme. Let $B$ be a quasi-compact and quasi-separated
algebraic space over $S$. Let $X \to B$ be a separated, finite type morphism.
Then $X$ has a compactification over $B$.
\end{theorem}

\begin{proof}
We first reduce to the Noetherian case. We strongly urge the reader
to skip this paragraph. First, we may replace $S$ by $\Spec(\mathbf{Z})$.
See Spaces, Section \ref{spaces-section-change-base-scheme} and
Properties of Spaces, Definition \ref{spaces-properties-definition-separated}.
There exists a closed immersion
$X \to X'$ with $X' \to B$ of finite presentation and separated.
See Limits of Spaces, Proposition
\ref{spaces-limits-proposition-separated-closed-in-finite-presentation}.
If we find a compactification of $X'$ over $B$, then
taking the scheme theoretic closure of $X$ in this will give
a compactification of $X$ over $B$. Thus we may assume
$X \to B$ is separated and of finite presentation.
We may write $B = \lim B_i$ as a directed
limit of a system of Noetherian algebraic spaces
of finite type over $\Spec(\mathbf{Z})$
with affine transition morphisms.
See Limits of Spaces, Proposition \ref{spaces-limits-proposition-approximate}.
We can choose an $i$ and a morphism $X_i \to B_i$ of finite
presentation whose base change to $B$ is $X \to B$, see
Limits of Spaces, Lemma \ref{spaces-limits-lemma-descend-finite-presentation}.
After increasing $i$ we may assume $X_i \to B_i$ is separated, see
Limits of Spaces, Lemma
\ref{spaces-limits-lemma-descend-separated-morphism}.
If we can find a compactification of $X_i$ over $B_i$, then the
base change of this to $B$ will be a compactification of $X$ over $B$.
This reduces us to the case discussed in the next paragraph.

\medskip\noindent
Assume $B$ is of finite type over $\mathbf{Z}$ in addition to being
quasi-compact and quasi-separated. Let $U \to X$ be an \'etale
morphism of algebraic spaces such that $U$ has a compactification
$Y$ over $\Spec(\mathbf{Z})$. The morphism
$$
U \longrightarrow B \times_{\Spec(\mathbf{Z})} Y
$$
is separated and quasi-finite by Morphisms of Spaces, Lemma
\ref{spaces-morphisms-lemma-monomorphism-loc-finite-type-loc-quasi-finite}
(the displayed morphism factors into an immersion hence is a monomorphism).
Hence by Zariski's main theorem (More on Morphisms of Spaces, Lemma
\ref{spaces-more-morphisms-lemma-quasi-finite-separated-pass-through-finite})
there is an open immersion of $U$ into an algebraic space $Y'$
finite over $B \times_{\Spec(\mathbf{Z})} Y$. Then $Y' \to B$ is proper
as the composition $Y' \to B \times_{\Spec(\mathbf{Z})} Y \to B$
of two proper morphisms
(use Morphisms of Spaces, Lemmas \ref{spaces-morphisms-lemma-finite-proper},
\ref{spaces-morphisms-lemma-composition-proper}, and
\ref{spaces-morphisms-lemma-base-change-proper}).
We conclude that $U$ has a compactification over $B$.

\medskip\noindent
There is a dense open subspace $U \subset X$ which is a scheme.
(Properties of Spaces, Proposition
\ref{spaces-properties-proposition-locally-quasi-separated-open-dense-scheme}).
In fact, we may choose $U$ to be an affine scheme
(Properties, Lemmas \ref{properties-lemma-Noetherian-irreducible-components}
and \ref{properties-lemma-maximal-points-affine}).
Thus $U$ has a compactification over $\Spec(\mathbf{Z})$;
this is easily shown directly but also follows from the
theorem for schemes, see
More on Flatness, Theorem \ref{flat-theorem-nagata}.
By the previous paragraph $U$ has a compactification over $B$. 
By Noetherian induction we can find a maximal dense open subspace
$U \subset X$ which has a compactification over $B$. We will show
that the assumption that $U \not = X$ leads to a contradiction.
Namely, by Lemma \ref{lemma-filter-Noetherian-space}
we can find a strictly larger open $U \subset W \subset X$
and a distinguished square $(U \subset W, f : V \to W)$
with $V$ affine and $U \times_W V$ dense image in $U$.
Since $V$ is affine, as before it has a compactification over $B$. Hence
Lemma \ref{lemma-two-compactifications}
applies to show that $W$ has a compactification over $B$
which is the desired contradiction.
\end{proof}
